\chapter{Background}
\label{sec:background}
\section{Phylogenetic Trees}
    Tree model and how to formulate likelihood on tree space 

\subsection{Evolution Model: Continuous-Time Markov Models}
    What does the (generated) data look like? 
        Generation using evolution models (specifically CTMCs) and transition diagrams 
\subsection{Tree Priors: Yule models}
    Tree priors (specifically Yule and others that I will use) 

\section{Bayesian Phylogenetic Inference}
\label{sec:bayesandmcmc}
\subsection{Likelihood Computation}
\label{sec:likelihoodcomputation}
\subsection{Posterior Sampling with Markov Chain Monte Carlo methods}
\subsection{Implementation in BEAST}

\section{Independent Paradigm}
        likelihood calculations, quality checks and convergence guarantees under such a model 
        
        can make this section + model misspecification its own chapter if its too long
\section{Dependent Paradigm}
\subsection{ZIP Construction}
\subsection{CTMC Product Construction}
Transition diagram
\subsection{Likelihood Computation}

\section{Model Misspecification}
    Model misspecification induced by dependence
        Data generation under P\_dependent but being evaluated on model under P\_independent

\section{Evaluation Metrics}
\subsection{Tree Recovery Metrics}
        Robinson-Foulds (RF) and nRF 
        Clade recovery (over 90\% of trees in posterior) 
        ... and maybe more if time allows 
\subsection{MCMC Diagnostics}
        Autocorrelation, and thus ESS 
        Mixing of chains (visual) 
        Thinning of data perhaps 

\section{ (TBD) helpful things (subequations, algorithms)}
To write equations over multiple lines (like systems of equation or equations too long to fit the page), one can use the \textit{align} environment coupled to the \textit{subequations} environment like so
\begin{subequations}
\begin{align}
    \mathbb{P}(0 \leadsto -a) &= D\int_0^\infty dt\, k_r e^{-k_r t}\int_0^t d\tau\,\partial_x u \big|_{x=-a}, \\
    \mathbb{P}(0 \leadsto b) &= -D\int_0^\infty dt\, k_r e^{-k_r t}\int_0^t d\tau\,\partial_x u \big|_{x=b}.
\end{align}
\end{subequations}

\begin{algorithm}[]
\SetAlgoLined
\SetKwInOut{KwInput}{Input}
\SetKwInOut{KwOutput}{Output}
\SetKwInOut{KwPre}{Pre}
\SetKw{Return}{return}
\SetKwProg{Fn}{Function}{}{end}
\LinesNumbered
\KwInput{$\textbf{m}$, such that $\mathbf{m}_i$ is the position of the $i$'th monitor\newline
$\textbf{l}$, such that $\mathbf{l}_i$ is the position of the $i$'th landmark\newline
$\mathbf{p}^m$, such that $\mathbf{p}^m_i$ is the ping latency from monitor $i$ to the target\newline
$\mathbf{p}^l$, such that $\mathbf{p}^l_i$ is the set of ping latencies to landmark $i$}

\BlankLine
\KwPre{Compute $\hat{p}_m(d \mid l)$, an estimator giving the likelihood of the target being distance $d$ away from the monitor, given that the monitor records a latency of $l$ to that target. Implemented by training a KDE using $\mathbf{p}^l$.\newline
Compute $\hat{p}_l(d \mid v)$, an estimator giving the likelihood of the target being distance $d$ away from the landmark, given a Canberra distance of $v$ between the target and the landmark, using training targets.
}
\BlankLine
\KwOutput{Most likely location of the target}
\BlankLine

\Fn{Likelihood($x$, $\mathbf{v}$)} {
MonitorScore $\gets \sum_{i=0}^{|\mathbf{m}|} \log{\hat{p}_m(d(x, \mathbf{m}_i) \mid l_i)}$\;
LandmarkScore $\gets \sum_{i=0}^{|\mathbf{l}|} \log{\hat{p}_l(d(x, \mathbf{l}_i) \mid \mathbf{v}_i)}$\;
\Return MonitorScore + LandmarkScore
}

\BlankLine
$\mathbf{v} \gets $\{$\mathrm{canberra\_distance}(\mathbf{l}_i, \mathbf{p}^m) \mid \mathbf{l}_i \in \mathbf{l}$\}

$\mathbf{C}$ $\gets$ Constraint-Based-Geolocation($\mathbf{m}$, $\mathbf{p}^m$)\;
$\mathbf{C_l}$ $\gets$ \{$m \in \mathbf{m} \mid \mathbf{C}$ contains $m\} \cup \{l \in \mathbf{l} \mid \mathbf{C}$ contains $l$\}\;
\BlankLine
\Return argmax$_{x\in \mathbf{C_l}}$ Likelihood($x$)

 \caption{Algorithm example}
 \label{algorithm:posit}
\end{algorithm}
